% ===========================================================================
\section{The method}\label{sec:method}
\warmth{0}\quad\whisper{The gradient points uphill. Negate it, take a step, repeat.}

\begin{strand}
At a point $x$, the gradient $\nabla f(x)$ is the direction of \emph{steepest ascent}, and its
negative the steepest descent. So the greedy thing to do is step a little way against the gradient
and reassess. The only knob is the \keyword{step size} $\eta$: too small and you crawl, too large
and you overshoot the valley and may climb the far wall.
\end{strand}

\block{The update}

\begin{definition}[gradient descent]\label{def:gd}
Starting from $x_0$, iterate
\[
  x_{k+1}\;=\;x_k-\eta\,\fillin[2cm],
\]
with step size $\eta>0$. The fixed points are exactly the \keyword{stationary points}
$\nabla f=0$; for convex $f$ those are the global minimizers.
\end{definition}

\recall{Why $-\nabla f$ and not some other descent direction? Among all unit directions $d$, the first-order change $\inner{\nabla f(x)}{d}$ is most negative when $d\propto-\nabla f(x)$. Steepest descent is gradient descent.}

\begin{yourturn}
Test it on a one-dimensional bowl $f(x)=\tfrac12\,a\,x^2$ with $a>0$.
\begin{itemize}[leftmargin=1.3em,itemsep=1pt]
  \item The update becomes $x_{k+1}=(1-\eta a)\,x_k$, so $x_k=\fillin[2.5cm]\cdot x_0$.
  \item It converges to $0$ exactly when $\abs{1-\eta a}<1$, i.e.\ for $0<\eta<\fillin[1.5cm]$.
        The sweet spot $\eta=1/a$ lands in \emph{one} step. \recall{$x_k=(1-\eta a)^k x_0$; converges iff $0<\eta<2/a$; $\eta=1/a$ is exact here because the curvature is constant.}
\end{itemize}
\workspace[2]
\end{yourturn}

\begin{strand}
That toy already shows the whole story. The good step size is tied to the \keyword{curvature} $a$:
match it and you fly, exceed $2/a$ and you diverge. In many dimensions there are \emph{many}
curvatures at once (the eigenvalues of the Hessian), and a single $\eta$ must serve all of them.
The spread of those curvatures is the tax of the next two sections.
\end{strand}

\loose{If $f$ is not differentiable (a kink, an $\ell_1$ penalty), replace $\nabla f$ by a \keyword{subgradient}; the method still works, just more slowly, and the proximal variants of §\ref{sec:faces} handle the kink exactly.}
